\setcounter{section}{5}

\Needspace{5\baselineskip}
\hypertarget{ppo-in-reinforcement-fine-tuning}{%
\section{PPO in Reinforcement Fine-Tuning}\label{ppo-in-reinforcement-fine-tuning}}

\Needspace{5\baselineskip}
\hypertarget{i.-ppos-objective-the-same-as-in-the-reinforcement-learning-part}{%
\subsection{I. PPO's Objective (the Same as in the Reinforcement Learning Part)}\label{i.-ppos-objective-the-same-as-in-the-reinforcement-learning-part}}

\Needspace{5\baselineskip}
Policy objective:

\[
J^{\mathrm{CLIP}}(\theta)
=
\mathbb{E}_t\left[
\min\left(
r_t(\theta)A_t,
\mathrm{clip}(r_t(\theta), 1-\epsilon, 1+\epsilon)A_t
\right)
\right]
\]

\Needspace{5\baselineskip}
Here:

\[
r_t(\theta)=\frac{\pi_\theta(a_t \mid s_t)}{\pi_{\theta_{\mathrm{old}}}(a_t \mid s_t)}
\]

\begin{itemize}
\tightlist
\item
  \(r_t(\theta)\): the probability ratio between the new and old policies for action \(a_t\).
\item
  \(A_t\): the estimated advantage of the current action, indicating how much better it is than average.
\item
  \(\epsilon\): the clipping range that limits policy-update magnitude.
\Needspace{5\baselineskip}
  Estimating \(A_t\):
\end{itemize}

\Needspace{5\baselineskip}
Looking only one step ahead, the \(Q\) value of action \(a_t\) can be estimated by \(r_t + \gamma V(s_{t+1})\). The estimate of \(A_t\) then equals the TD error:

\[
\hat{A}_t^{(1)}
=
\underbrace{r_t + \gamma V(s_{t+1})}_{\approx Q(s_t,a_t)}
- V(s_t)
=
\delta_t
\]

\Needspace{5\baselineskip}
PPO uses GAE by default to calculate advantage, taking an exponentially weighted moving average of TD errors:

\[
\hat{A}_t^{\mathrm{GAE}}
=
\delta_t
+
(\gamma\lambda)\delta_{t+1}
+
(\gamma\lambda)^2\delta_{t+2}
+
\cdots
\]

Here, \(\lambda\) is a hyperparameter in \([0,1]\).

\Needspace{5\baselineskip}
The effect of TD errors is controlled through \(\lambda\):

\begin{itemize}
\tightlist
\item
  At \(\lambda = 0\): high bias, low variance.
\end{itemize}

\[
\hat{A}_t = \delta_t = r_t + \gamma V(s_{t+1}) - V(s_t)
\]

\(A_t\) is exactly the current one-step TD error.

\begin{itemize}
\tightlist
\item
  Advantage: low variance, depending only on one random reward \(r_t\).
\item
  Disadvantage: high bias, strongly dependent on critic accuracy at \(V(s_{t+1})\). A poorly trained critic produces incorrect estimates.
\item
  At \(\lambda = 1\): no bias, high variance.
\end{itemize}

\[
\hat{A}_t
=
\sum_{l=0}^{\infty}\gamma^l\delta_{t+l}
=
\left(\sum_{l=0}^{\infty}\gamma^l r_{t+l}\right)
- V(s_t)
\]

Expanding the expression cancels the intermediate \(V\) terms, leaving Monte Carlo returns (actual returns) minus the baseline.

\begin{itemize}
\tightlist
\item
  Advantage: low bias; actual returns are the most accurate facts.
\item
  Disadvantage: very high variance from accumulating randomness at every environment step.
\end{itemize}

\Needspace{5\baselineskip}
\hypertarget{ii.-ppo-workflow-in-reinforcement-fine-tuning-slightly-different-from-the-earlier-version}{%
\subsection{II. PPO Workflow in Reinforcement Fine-Tuning (Slightly Different from the Earlier Version)}\label{ii.-ppo-workflow-in-reinforcement-fine-tuning-slightly-different-from-the-earlier-version}}

\Needspace{4\baselineskip}
\begin{enumerate}
\def\labelenumi{\arabic{enumi}.}
\tightlist
\item
  Data collection
\end{enumerate}

In traditional RL environments such as robot control, the agent interacts continuously for, say, 2,048 steps, producing a long trajectory that is later shuffled into minibatches for gradient calculation. PPO for LLM fine-tuning instead uses vectorized environments: multiple prompts can be processed in parallel, independently generating trajectories stored in a replay buffer.

``2,048 steps'' is actually the total number of tokens generated across all parallel environments.

Suppose GPU parallelism allows \(N\) prompts to be processed simultaneously (the number of parallel environments, such as \(N = 16\)).

\begin{itemize}
\tightlist
\item
  At \(t = 1\), the model sees 16 different prompts and generates 16 tokens concurrently, one per prompt. This counts as 16 steps.
\item
  At \(t = 2\), it generates another 16 tokens based on the preceding tokens.
\item
  The process continues.
\end{itemize}

\Needspace{5\baselineskip}
The calculation is:

\[
\text{Total Steps }(2048)
=
\text{Number of parallel prompts}
\times
\text{Average response length}
\]

\Needspace{5\baselineskip}
For example, ``run 2,048 steps'' in a diagram may mean:

\begin{enumerate}
\def\labelenumi{\arabic{enumi}.}
\tightlist
\item
  The system samples 16 different prompts from the dataset.
\item
  The model generates answers to all 16 in parallel.
\item
  Each answer ends after an average of 128 tokens.
\item
  Total collected steps: \(16 \times 128 = 2048\).
\end{enumerate}

Thus, 2,048 steps contain complete stories for 16 prompts, not one extremely long story for a single prompt.

\Needspace{4\baselineskip}
\begin{enumerate}
\def\labelenumi{\arabic{enumi}.}
\setcounter{enumi}{1}
\tightlist
\item
  Advantage calculation (the same as traditional RL)
\end{enumerate}

Initialization: initialize actor \(\pi_\theta\) and critic \(V_\phi\).

\Needspace{5\baselineskip}
For each iteration:

\begin{enumerate}
\def\labelenumi{\arabic{enumi}.}
\tightlist
\item
\Needspace{5\baselineskip}
  Data collection (interaction):

  \begin{itemize}
  \tightlist
  \item
    Run current policy \(\pi_{\theta_{\mathrm{old}}}\) for \(T\) environment steps, such as 2,048.
  \item
    Collect trajectories \(\{s_t, a_t, r_t, s_{t+1}, \log \pi_{\theta_{\mathrm{old}}}(a_t \mid s_t)\}\).
  \item
    Calculate state values \(V(s_t)\) with the critic.
  \end{itemize}
\item
\Needspace{5\baselineskip}
  Advantage calculation:

  \begin{itemize}
  \tightlist
  \item
    Calculate TD error \(\delta_t = r_t + \gamma V(s_{t+1}) - V(s_t)\).
  \item
    Calculate generalized advantage estimation (GAE) \(A_t\), a recursive formula balancing bias and variance.
  \end{itemize}
\item
  Optimization and policy synchronization
\end{enumerate}

Collection: like a harvester, the GPU processes many prompts in parallel to fill a replay buffer of capacity 2,048 rapidly. It may contain question--answer records for 16 different problems.

Shuffle: thoroughly shuffle the 2,048 tokens (data points) to stabilize training and break correlations.

Minibatches: split the shuffled data into small batches, such as batch size 64.

\begin{itemize}
\tightlist
\item
  This gives \(2048 / 64 = 32\) minibatches.
\item
  The model uses those 32 minibatches to update its parameters.
\end{itemize}

\Needspace{4\baselineskip}
\begin{enumerate}
\def\labelenumi{\arabic{enumi}.}
\setcounter{enumi}{2}
\tightlist
\item
\Needspace{5\baselineskip}
  Optimization:
\end{enumerate}

\begin{itemize}
\tightlist
\item
  Important: \(\pi_{\theta_{\mathrm{old}}}\) remains fixed as the denominator. The new \(\pi_\theta\) in the numerator is optimized.
\item
  Shuffle the \(T\) collected data points into minibatches, such as 64 points each.
\item
\Needspace{5\baselineskip}
  Repeat for multiple epochs, such as 10:

  \begin{enumerate}
  \def\labelenumi{\arabic{enumi}.}
  \tightlist
  \item
    For each minibatch, calculate the new probability ratio \(r_t(\theta)=\frac{\pi_\theta(a \mid s)}{\pi_{\theta_{\mathrm{old}}}(a \mid s)}\).
  \item
    Calculate the clipping loss \(L^{\mathrm{CLIP}}\).
  \item
    Calculate critic value loss \(L^{\mathrm{VF}} = (V_\phi(s_t) - V_{\mathrm{target}})^2\).
  \item
    Calculate entropy regularization \(S\) to encourage exploration.
  \item
    Total loss: \(L = -L^{\mathrm{CLIP}} + c_1 L^{\mathrm{VF}} - c_2 S\). Note the signs for gradient ascent/descent.
  \item
    Backpropagate to update \(\theta\) and \(\phi\).
  \end{enumerate}
\end{itemize}

\Needspace{4\baselineskip}
\begin{enumerate}
\def\labelenumi{\arabic{enumi}.}
\setcounter{enumi}{3}
\tightlist
\item
\Needspace{5\baselineskip}
  Policy synchronization:
\end{enumerate}

\begin{itemize}
\tightlist
\item
  After the round of updates, set \(\pi_{\theta_{\mathrm{old}}} \leftarrow \pi_\theta\).
\item
  Clear the data buffer and begin the next iteration.
\end{itemize}

\Needspace{5\baselineskip}
\hypertarget{iii.-why-do-llms-choose-ppo-over-sac-and-on-policy-over-off-policy}{%
\subsection{III. Why Do LLMs Choose PPO over SAC, and On-Policy over Off-Policy?}\label{iii.-why-do-llms-choose-ppo-over-sac-and-on-policy-over-off-policy}}

\Needspace{5\baselineskip}
\hypertarget{llms-output-discrete-rather-than-continuous-actions}{%
\subsubsection{(1) LLMs Output Discrete Rather than Continuous Actions}\label{llms-output-discrete-rather-than-continuous-actions}}

SAC's original design: SAC is a maximum-entropy RL algorithm designed for continuous action spaces. In robot control, actions are real-valued joint angles or torques. Discrete SAC variants exist, but high-dimensional discrete spaces are extremely expensive to handle.

The LLM setting: generation selects the next token from a huge vocabulary.

\begin{itemize}
\tightlist
\item
  Action-space size \(|A|\): usually \(32,000\) to \(128,000\) or larger, such as the vocabularies of GPT-4 or Qwen.
\item
  Curse of dimensionality: an actor--critic's critic evaluates \(Q(s,a)\).

  \begin{itemize}
  \tightlist
  \item
    Discrete SAC generally marginalizes over all actions or calculates Softmax to obtain target \(Q\) values and policy entropy.
  \item
\Needspace{5\baselineskip}
    Its formulas involve the expectation of \(\pi(a' \mid s') \cdot [Q(s',a') - \alpha \log \pi(a' \mid s')]\), namely:
  \end{itemize}
\end{itemize}

\[
\mathbb{E}_{a'\sim \pi(\cdot \mid s')}
\left[
Q(s',a')-\alpha\log\pi(a'\mid s')
\right]
=
\sum_{a'\in A}\pi(a'\mid s')
\left[
Q(s',a')-\alpha\log\pi(a'\mid s')
\right]
\]

Each update therefore requires a full probability calculation over \(50,000+\) tokens, which is extremely costly.

One might suggest that a 100,000-token vocabulary already resembles a continuous space. Would changing the policy output from a 100,000-dimensional probability distribution to a continuous embedding vector not greatly reduce memory and compute? In practice, this approach has so far proved unworkable.

\Needspace{4\baselineskip}
\begin{enumerate}
\def\labelenumi{\arabic{enumi}.}
\tightlist
\item
  Blocked gradients: differentiating across a discrete cliff
\end{enumerate}

This is the central mathematical obstacle.

\begin{itemize}
\tightlist
\item
  The illusion of continuous actions: suppose the actor outputs a continuous vector \(\mathbf{v}\in\mathbb{R}^d\), such as \(d=4096\). To provide it to the environment or predict the next token, it must be mapped back to the real vocabulary.
\item
\Needspace{5\baselineskip}
  Nearest-neighbor search: find the token embedding \(\mathbf{e}_k\) closest to \(\mathbf{v}\) in the embedding table:
\end{itemize}

\[
k
=
\underset{j\in V}{\arg\min}\,
\lVert \mathbf{v}-\mathbf{e}_j\rVert_2
\]

\begin{itemize}
\tightlist
\item
  Gradient difficulty: \(\arg\min\) is nondifferentiable, or has zero gradient almost everywhere.

  \begin{itemize}
  \tightlist
  \item
    If \(\mathbf{v}\) moves slightly without crossing a nearest-neighbor boundary, the nearest token remains \(k\). Output and loss do not change, so the gradient is 0.
  \item
    Crossing a boundary changes the nearest neighbor to token \(m\), creating a step discontinuity with an infinite or undefined derivative.
  \end{itemize}
\end{itemize}

Without gradients, backpropagation stops and the actor learns nothing. Relaxations such as Gumbel-Softmax exist, but have enormous variance and unstable training in a 100,000-dimensional space.

\Needspace{4\baselineskip}
\begin{enumerate}
\def\labelenumi{\arabic{enumi}.}
\setcounter{enumi}{1}
\tightlist
\item
  Empty regions: sparsity of the embedding space
\end{enumerate}

One might ask, ``Why use argmax at all? Why not feed continuous output \(\mathbf{v}\) directly into the next network layer?''

This involves the manifold hypothesis.

\begin{itemize}
\tightlist
\item
  Word vectors are not uniformly distributed: in 4,096-dimensional space, valid token embeddings occupy only a tiny point cloud. Most of the space is invalid void space.
\item
  Continuous-control difficulty: RL exploration usually adds Gaussian noise to actions. Adding noise to \(\mathbf{v}\) will probably place it in an empty region.

  \begin{itemize}
  \tightlist
  \item
    Result: the vector resembles neither ``apple'' nor ``banana,'' but is disordered and semantically meaningless.
  \item
    Language-model collapse: passing this meaningless vector onward makes Transformer attention process confused information, causing subsequent generation to become gibberish.
  \end{itemize}
\end{itemize}

\Needspace{4\baselineskip}
\begin{enumerate}
\def\labelenumi{\arabic{enumi}.}
\setcounter{enumi}{2}
\tightlist
\item
  The efficiency paradox: KNN is slower than matrix multiplication
\end{enumerate}

The continuous approach aims to improve efficiency, but may be slower.

\begin{itemize}
\tightlist
\item
  Discrete approach (current practice): the output layer is a linear transformation \(W\cdot h_o\). Although \(W\) is large, for example \(4096\times100000\), it is just a large matrix multiplication, heavily optimized on GPUs through Tensor Cores.
\item
  Continuous approach (proposed): converting the actor's output vector into a word for human readers or reward computation requires searching the entire vocabulary for the nearest neighbor.

  \begin{itemize}
  \tightlist
  \item
    Each generation step calculates distances to 100,000 embeddings.
  \item
    This still amounts to multiplying \(1\times4096\) by \(4096\times100000\), plus sorting or \(\arg\min\).
  \item
    It is no faster than directly outputting logits and is slower because of additional indexing.
  \end{itemize}
\end{itemize}

\Needspace{5\baselineskip}
\hypertarget{llms-are-better-suited-to-on-policy-than-off-policy-learning}{%
\subsubsection{(2) LLMs Are Better Suited to On-Policy than Off-Policy Learning}\label{llms-are-better-suited-to-on-policy-than-off-policy-learning}}

In traditional RL, such as robot control and game AI, much effort goes into off-policy learning because sample efficiency is considered essential.

For LLMs, on-policy learning is not merely a compromise; it can be a core advantage, overturning traditional RL rules of thumb. This can be analyzed mathematically, computationally, and through alignment objectives.

\Needspace{4\baselineskip}
\begin{enumerate}
\def\labelenumi{\arabic{enumi}.}
\tightlist
\item
  Mathematical breakdown: the curse of importance sampling
\end{enumerate}

\Needspace{5\baselineskip}
Off-policy learning uses data from an old policy to update a new one. Correcting distribution differences generally requires importance-sampling (IS) weights:

\[
\rho_t
=
\frac{\pi_{\mathrm{new}}(a_t \mid s_t)}
{\pi_{\mathrm{old}}(a_t \mid s_t)}
\]

\Needspace{5\baselineskip}
In robot control, action sequences are usually short or action spaces simple, keeping \(\rho_t\) manageable. In LLMs, the situation becomes uncontrolled:

\begin{itemize}
\tightlist
\item
\Needspace{5\baselineskip}
  Multiplication across sequence length: a trajectory usually contains hundreds or thousands of tokens. Its probability ratio is the product of per-step ratios:
\end{itemize}

\[
\rho_{\mathrm{trajectory}}
=
\prod_{t=1}^{T}
\frac{\pi_{\mathrm{new}}(a_t\mid s_t)}
{\pi_{\mathrm{old}}(a_t\mid s_t)}
\]

\Needspace{5\baselineskip}
Off-policy algorithms update new policies using old-policy data. Depending on value estimation, there are two cases:

\Needspace{4\baselineskip}
\begin{enumerate}
\def\labelenumi{(\arabic{enumi})}
\tightlist
\item
  Without a Q network
\end{enumerate}

\Needspace{5\baselineskip}
The new-policy objective is:

\[
J(\pi_{\mathrm{new}})
=
\mathbb{E}_{\tau\sim\pi_{\mathrm{new}}}[R(\tau)]
=
\int P(\tau\mid\pi_{\mathrm{new}})R(\tau)d\tau
\]

\Needspace{5\baselineskip}
Since only \(\pi_{\mathrm{old}}\) data is available, importance sampling is required:

\[
J(\pi_{\mathrm{new}})
=
\mathbb{E}_{\tau\sim\pi_{\mathrm{old}}}
\left[
\frac{P(\tau\mid\pi_{\mathrm{new}})}
{P(\tau\mid\pi_{\mathrm{old}})}
R(\tau)
\right]
\]

\Needspace{5\baselineskip}
Expanding trajectory probability \(P(\tau)\) cancels the environment dynamics, leaving action-probability ratios:

\[
\frac{P(\tau\mid\pi_{\mathrm{new}})}
{P(\tau\mid\pi_{\mathrm{old}})}
=
\frac{\prod_{t=0}^{T}\pi_{\mathrm{new}}(a_t\mid s_t)}
{\prod_{t=0}^{T}\pi_{\mathrm{old}}(a_t\mid s_t)}
=
\prod_{t=0}^{T}
\frac{\pi_{\mathrm{new}}(a_t\mid s_t)}
{\pi_{\mathrm{old}}(a_t\mid s_t)}
=
\prod_{t=0}^{T}\rho_t
\]

\Needspace{5\baselineskip}
This is where the product arises:

\[
\text{Off-Policy Gradient}
\propto
\mathbb{E}
\left[
\left(\prod_{t=0}^{T}\rho_t\right)R(\tau)
\right]
\]

\Needspace{5\baselineskip}
For LLMs, \(T\approx1024\). Even small deviations of \(\rho_t\) from 1 make the product fluctuate sharply between 0 and \(+\infty\), exploding variance. With \(T=1024\), even small per-step differences give:

\begin{itemize}
\tightlist
\item
  Ratio \(1.01\): \(1.01^{1024}\approx 26,739\).
\item
  Ratio \(0.99\): \(0.99^{1024}\approx 0.00003\).
\end{itemize}

IS-weight variance becomes nearly uncontrollable. PPO's on-policy advantage is that it enforces \(\pi_{\mathrm{new}}\approx\pi_{\mathrm{old}}\), limits policy drift through clipping and short rollouts, and keeps \(\rho_t\approx1\), avoiding long-sequence numerical explosions.

\Needspace{4\baselineskip}
\begin{enumerate}
\def\labelenumi{(\arabic{enumi})}
\setcounter{enumi}{1}
\tightlist
\item
  With a Q network: avoiding exploding variance but risking exploding bias
\end{enumerate}

Generating an LLM response usually takes \(T=1000\) steps. Suppose the Q network has a small estimation error \(\epsilon\), nearly inevitable with a 70B-parameter model and sparse rewards.

\[
Q_{\mathrm{target}}(s_{T-1})
=
r_{T-1}
+\mathrm{EndValue}
+\epsilon
\]

\[
Q_{\mathrm{target}}(s_{T-2})
=
r_{T-2}
+\gamma Q_{\mathrm{target}}(s_{T-1})
+\epsilon
\]

\Needspace{5\baselineskip}
Recursing backward accumulates error along the long sequence:

\[
Q_{\mathrm{target}}(s_0)
=
\sum_{t=0}^{T-1}\gamma^t r_t
+
\sum_{k=0}^{T-1}\gamma^k \epsilon
\]

Because of Q-value overestimation, \(\epsilon\) is often positive, and its accumulation undermines critic judgments.

\Needspace{4\baselineskip}
\begin{enumerate}
\def\labelenumi{\arabic{enumi}.}
\setcounter{enumi}{1}
\tightlist
\item
  Reversed computational architecture: recomputation costs more than generation
\end{enumerate}

In traditional RL, environmental interaction takes physical time while neural computation is fast, so stored data should be reused repeatedly on GPUs.

In LLM RLHF, interaction (text generation) is itself neural computation and is extremely expensive.

Using old data with off-policy methods such as SAC or IMPALA requires knowing its probability distribution under current parameters to compute \(\log\pi_{\mathrm{new}}\) or KL divergence. This requires another forward pass of the old prompt and response through the current 70B model.

\Needspace{4\baselineskip}
\begin{enumerate}
\def\labelenumi{\arabic{enumi}.}
\setcounter{enumi}{2}
\tightlist
\item
  Distinctive alignment objectives: fine-tuning rather than seeking a global optimum
\end{enumerate}

This is the most fundamental difference.

\begin{itemize}
\tightlist
\item
  Traditional RL, such as AlphaGo: seek a global optimum, discovering brilliant moves humans have never seen. Off-policy learning retains accidentally discovered unusual high-reward actions in a replay buffer and repeatedly reinforces them.
\item
  LLM alignment (RLHF): constrained fine-tuning. The goal is not to invent a new language or hack rewards through gibberish, but to remain near the SFT model with small style and safety improvements.

  \begin{itemize}
  \tightlist
  \item
    Off-policy risk: exploring regions far from the current policy can cause model collapse, in which the model speaks nonsense to maximize reward.
  \end{itemize}
\end{itemize}

Off-policy learning's tendency to deviate comes from its own characteristics.

Mechanism: store all historical transitions \((s,a,r,s')\) in a replay buffer and sample them randomly during training.

Consequence: the buffer may contain data from \(\pi_{\mathrm{old}}\) 100 epochs earlier.

\begin{itemize}
\tightlist
\item
  That policy may have been highly random, accidentally trying an extreme action early in exploration. In an LLM, this could be an obscure gibberish word mistakenly given a high score by the reward model.
\item
  Although mature \(\pi_{\mathrm{new}}\) would never intentionally produce it, off-policy sampling retrieves it and notices, ``This action scored 100!''
\item
  The update then pulls the current policy toward that historical, unreasonable high-reward region.
\end{itemize}

By contrast, on-policy PPO discards data after use. If an unusual high-scoring error occurs once and does not recur, PPO quickly forgets it rather than repeatedly revisiting it.

On-policy learning calculates the policy gradient under the current policy and takes one step to adjust the distribution, generally staying near the existing solution. The ``Phenomenon Explanation and Theoretical Analysis in Large-Model Research'' part also proves that on-policy methods repeatedly project between feasible and optimal policies. The limiting policy is theoretically the optimal policy with the smallest KL divergence from the initial policy.

Overestimating unseen regions is another reason for deviation.

Q-value generalization error: neural networks generalize, so a Q network predicts values even for unseen state--action pairs \((s,a_{\mathrm{unseen}})\) that current policy \(\pi\) has never visited.

Overestimation tendency: without real-data constraints on these \((s,a_{\mathrm{unseen}})\) pairs, the Q network's predictions often have high variance. The \(\max\) in \(y=r+\gamma\max Q\) specifically selects incorrectly overestimated values.

Result: the Q network effectively imagines, ``I have never been there, but it must be full of gold.''

Action: the actor is drawn toward regions far from the current distribution that may not actually yield high returns.

\Needspace{5\baselineskip}
\hypertarget{llm-fine-tuning-requires-controlled-update-sizes}{%
\subsubsection{(3) LLM Fine-Tuning Requires Controlled Update Sizes}\label{llm-fine-tuning-requires-controlled-update-sizes}}

LLM fine-tuning must control update magnitudes. Excessively large parameter updates rapidly cause ``language collapse'': maximizing reward destroys pretrained language ability, yielding gibberish or repeated meaningless characters, another form of reward hacking.

\Needspace{5\baselineskip}
PPO limits the difference between new and old policies through clipping:

\[
L^{\mathrm{CLIP}}(\theta)
=
\mathbb{E}_t
\left[
\min
\left(
r_t(\theta)\hat{A}_t,
\mathrm{clip}(r_t(\theta),1-\epsilon,1+\epsilon)\hat{A}_t
\right)
\right]
\]

Here, \(r_t(\theta)=\frac{\pi_\theta(a_t\mid s_t)}{\pi_{\theta_{\mathrm{old}}}(a_t\mid s_t)}\) is the probability ratio.

\begin{itemize}
\tightlist
\item
  Role: \(\epsilon\), usually 0.2, prevents \(\pi_\theta\) from moving too far from \(\pi_{\theta_{\mathrm{old}}}\).
\item
  Result: while learning human preferences, the LLM retains pretrained abilities such as grammar and logic.
\end{itemize}

Besides limiting one update through clipping, modern large models often add a KL term relative to a reference distribution (the pretrained model's output distribution), preventing many small updates from accumulating excessive drift.

\Needspace{5\baselineskip}
SAC maximizes entropy alongside cumulative reward:

\[
J(\pi)
=
\sum_{t=0}^{T}
\mathbb{E}_{(s_t,a_t)\sim\rho_\pi}
\left[
r(s_t,a_t)
+\alpha\mathcal{H}(\pi(\cdot\mid s_t))
\right]
\]

Risk: although entropy encourages exploration, RLHF usually seeks confident, logically rigorous answers rather than high-entropy random tokens for exploration. Excessive entropy regularization can make LLM output divergent and incoherent.

\Needspace{5\baselineskip}
\hypertarget{q-network-stability-under-sparse-rewards}{%
\subsubsection{(4) Q-Network Stability Under Sparse Rewards}\label{q-network-stability-under-sparse-rewards}}

\begin{itemize}
\tightlist
\item
  SAC (tightly coupled actor and critic): SAC strongly depends on Q-network critic accuracy. Training an LLM critic to judge whether a sentence is good, effectively a proxy for the reward model, is already difficult. A divergent or inaccurate critic quickly corrupts the actor.

  \begin{itemize}
  \tightlist
  \item
    Under large-scale sparse rewards, bootstrapping in Q-learning methods easily produces overestimation bias.
  \end{itemize}
\item
  PPO (Monte Carlo-like regression): PPO uses generalized advantage estimation (GAE) to calculate \(\hat{A}_t\). Although it also uses a critic (value network), it relies more on actual returns over the whole trajectory. Compared with SAC's strongly bootstrapped TD learning, PPO is usually more robust for long-sequence generation with long episodes.
\end{itemize}

\FloatBarrier
\Needspace{5\baselineskip}
\hypertarget{references-84}{%
\subsection{References}\label{references-84}}

\vspace{0.25em}
\begin{itemize}
\tightlist
\item
  Schulman, J., Moritz, P., Levine, S., Jordan, M., \& Abbeel, P. (2015). \href{https://arxiv.org/abs/1506.02438}{High-Dimensional Continuous Control Using Generalized Advantage Estimation}. arXiv:1506.02438.
\item
  Schulman, J., Wolski, F., Dhariwal, P., Radford, A., \& Klimov, O. (2017). \href{https://arxiv.org/abs/1707.06347}{Proximal Policy Optimization Algorithms}. arXiv:1707.06347.
\item
  Haarnoja, T., Zhou, A., Abbeel, P., \& Levine, S. (2018). \href{https://arxiv.org/abs/1801.01290}{Soft Actor-Critic: Off-Policy Maximum Entropy Deep Reinforcement Learning with a Stochastic Actor}. ICML 2018.
\item
  Haarnoja, T., Zhou, A., Hartikainen, K., et al.~(2018). \href{https://arxiv.org/abs/1812.05905}{Soft Actor-Critic Algorithms and Applications}. arXiv:1812.05905.
\end{itemize}

\clearpage
